\begin{abstract}
This article merges two hundred and nine research reports that start from one input: the Riemann
hypothesis for the Riemann zeta function and for Dirichlet $L$-functions, as stated in the
Lean development \repo. The input is taken as given and is not re-verified here. The
reports ask two questions. What does the theorem settle elsewhere? And which methods inside
its proof carry over to other problems? Each result is stated once, in the strongest form
that the reports prove, and results that follow from stronger ones are left out.

The first group of results consists of consequences of the input: complete lists of
imaginary quadratic fields of class-group exponent six and seven, a classification of abelian
fields in a question on Mertens errors, statements on zeros, Euler products and prime races,
covariance and rigidity of smoothed prime errors, the recovery of zeros, moduli and
characters from prime data, and a range of arithmetic applications through published
conditional theorems. The second group transfers the one-sided analytic step of the proof.
It develops that step into criteria on sparse sets of scales and of twists, and into
sampling, recovery and positivity statements for Li coefficients. It reduces the Riemann hypothesis
for Dedekind zeta functions of nonabelian fields, and for several Artin, Hecke and modular
$L$-functions, to explicit one-sided inequalities over rational primes. These are
reductions: the inequalities are open, and no new Riemann hypothesis is proved. The third
group is unconditional. It contains sharp theorems on torsion specialization of $p$-adic
power series, on Lubin--Tate branches, on sparse equations for singular moduli, on sets
without geometric progressions in the integers, in number fields and in polynomial rings, and
on spanning-tree valuations in graph towers. A last part records what the arithmetic engine
of the proof needs, and where it stops, when one leaves Dirichlet characters.
\end{abstract}
